% -*- mode: LaTeX; -*-
%
% Alex Han -- curriculum vitae.
%
% Build:  latexmk -pdf AlexHan.tex        (clean up with: latexmk -C)
%
% This is the whole CV: preamble, then content, in one file. The only other
% file that matters is res.cls, the document class that puts section titles
% out in the left margin. res.cls is NOT part of TeX Live and cannot be
% reinstalled with tlmgr -- the copy in this repo is the only one.
%
%%%%%%%%%%%%%%%%%%%%%%%%%%%%%%%%%%%%%%%%%%%%%%%%%%%%%%%%%%%%%%%%%%%%%%%%%%%%%%
% ENTRY FORMATTING CHEAT SHEET
%
% There is no structured entry macro in this template -- entries are formatted
% by hand. Three idioms cover everything:
%
%   1. A standard entry: bold role/degree, plain institution, bold date at the
%      right margin. Note the trailing "\\" then a lone "\\" to separate it from
%      the next entry:
%
%        {\bf M.D.-Ph.D. Candidate,} Institution, City, ST \hfill {\bf 2025--Present}\\
%        \\
%
%   2. Sub-details underneath an entry, with the bullet suppressed by \item[ ]:
%
%        \begin{itemize}
%        \item[ ] Advisor: Prof.\ Name
%        \end{itemize}
%
%   3. A section whose first content is a list needs \secstartswithlist{} right
%      after \section{} to cancel a stray blank line.
%
% Dates use en-dashes: 2020--2024. Escape "&" as "\&" and "%" as "\%".
%
% WHEN YOU ADD A PUBLICATION: bump \setcounter{numPubs} below the Publications
% heading to the new total, or the descending numbering will not end at 1.
%%%%%%%%%%%%%%%%%%%%%%%%%%%%%%%%%%%%%%%%%%%%%%%%%%%%%%%%%%%%%%%%%%%%%%%%%%%%%%

\let\nofiles\relax % This is because res says to not emit aux files,
                   % but lastpage needs aux files.
\documentclass[margin,line]{res}
\usepackage[colorlinks=true]{hyperref}
\usepackage[utf8]{inputenc}
\usepackage[T1]{fontenc}
\usepackage{lmodern} % Scalable Type 1 fonts; without this, T1 + Computer
                     % Modern makes pdfTeX emit Type 3 bitmap fonts.
\usepackage{microtype}
\usepackage{amsmath,amssymb}
\usepackage{lastpage}
\usepackage{fancyhdr}
\usepackage{etaremune}
\usepackage[normalem]{ulem}
\usepackage[style=iso]{datetime2}

\oddsidemargin -.5in
\evensidemargin -.5in
\voffset -25pt
%\topmargin -.2in
\headsep 25pt
\textwidth=6.0in
\textheight=8.9in
\itemsep=0in
\parsep=0in

% Headings
\pagestyle{fancy}
\lhead{Alex Han --- Curriculum Vitae (as of \today)}
\chead{}
\rhead{\thepage\ of \pageref*{LastPage}}
\lfoot{}
\cfoot{}
\rfoot{}
\renewcommand{\headrulewidth}{0.4pt}
%\renewcommand{\footrulewidth}{0.4pt}

% Give hyperref some metadata
\hypersetup{pdftitle={Alex Han — Curriculum Vitae (CV)},
  pdfauthor={Alex Han},
  pdfsubject={Alex Han's academic curriculum vitae (CV)},
  pdfkeywords={medicine, MD-PhD, computational biology, human genetics,
    statistical genetics, machine learning, neurodegeneration,
    Alzheimer's disease, drug target discovery}}

% I want small caps for the section style
\def\sectionfont{\sc}

% Generate a PDF TOC
\let\oldsection\section
\def\section#1{\oldsection{#1}%
\phantomsection%
\addcontentsline{toc}{section}{#1}%
}

% The above seems to screw up the spacing for sections that start with lists...
\newcommand{\secstartswithlist}{\leavevmode \vspace{-\baselineskip}}

% Publication link helpers -- use whichever applies to a given entry.
\newcommand{\doi}[1]{[\href{https://doi.org/#1}{doi:#1}]}
\newcommand{\biorxiv}[1]{[\href{https://doi.org/#1}{bioRxiv:#1}]}
\newcommand{\medrxiv}[1]{[\href{https://doi.org/#1}{medRxiv:#1}]}
\newcommand{\pmid}[1]{[\href{https://pubmed.ncbi.nlm.nih.gov/#1/}{PMID:#1}]}

% Publications are listed newest-first but numbered DESCENDING (via etaremune),
% so the top of the list always shows your current total and #1 is your oldest.
\newcounter{numPubs}
\newcounter{pubCounter}

\begin{document}

\newcommand{\myname}{Alex Han}
\newlength{\mynamewidth}
\settowidth{\mynamewidth}{\namefont\myname}

\name{\hspace*{0.5\textwidth}\hspace{-0.5\mynamewidth} \myname \vspace*{.1in}}
% On the first page, have no header.
\thispagestyle{empty}

\begin{resume}

\section{Contact Information}
% Three lines, left column plain and right column hyperlinked. Keep the two
% columns balanced -- add or remove a line from both sides together.
%
% This repo is public. Street address and phone are left out on purpose;
% add them only if you want them on the open internet.
%
% TODO: fill in your real details.
Department / Program        \hfill \href{mailto:TODO@example.edu}{TODO@example.edu}\\
Institution                 \hfill \href{https://example.com/}{example.com}\\
City, ST 00000 USA          \hfill \href{https://orcid.org/0000-0000-0000-0000}{ORCID: 0000-0000-0000-0000}
% Optional extra line:
% 000 Building Name           \hfill \href{tel:1-000-000-0000}{1-000-000-0000}\\

\section{Education}
% TODO: replace with your degrees, most recent first.
{\bf M.D.-Ph.D. Candidate,} Institution, City, ST, USA \hfill {\bf 2025--Present}\\
\vspace*{-.1in}
\begin{itemize}
\item[ ] Program: TODO (e.g.\ Medical Scientist Training Program)
\item[ ] Advisor: TODO
\item[ ] Thesis Title: {\it TODO}
\end{itemize}

{\bf B.A./B.S., Field,} Institution, City, ST, USA \hfill {\bf 2020--2024}\\
\vspace*{-.1in}
\begin{itemize}
\item[ ] TODO: honors, distinctions, GPA if you want it
\item[ ] Senior Thesis Advisor: TODO
\end{itemize}

\section{Research Experience}
% TODO: one block per lab/position, most recent first.
{\bf Position,} Lab / Group, Institution, City, ST \hfill {\bf 2024--Present}\\
\vspace*{-.1in}
\begin{itemize}
\item[ ] TODO: one line on the question and your contribution
\item[ ] TODO: methods, tools, or datasets
\end{itemize}

{\bf Position,} Lab / Group, Institution, City, ST \hfill {\bf 2023--2025}\\
\vspace*{-.1in}
\begin{itemize}
\item[ ] TODO
\end{itemize}

% Author-name convention: bold your own name so it is findable at a glance,
% e.g. {\bf Han,~A.~L.}  Use ~ between initials to prevent line breaks.
\setcounter{numPubs}{3} % TODO: update whenever you add a publication
\setcounter{pubCounter}{\value{numPubs}}

\section{Manuscripts Under Review and Preprints}
\secstartswithlist{}%
\begin{etaremune}[start=\value{pubCounter}]
\item
  Lastname,~A.~B.,
  {\bf Han,~A.~L.},
  Lastname,~C.~D.
  (2026)
  {\it Title of the manuscript under review},
  submitted to Journal Name.
  \biorxiv{10.1101/0000.00.00.000000}
% TODO: add further preprints here
\end{etaremune}
% Carry the numbering into the next list, continuing one below the last item.
\setcounter{pubCounter}{\value{enumi}}\addtocounter{pubCounter}{-1}

\section{Peer-Reviewed Publications}
\secstartswithlist{}%
\begin{etaremune}[start=\value{pubCounter}]
\item
  {\bf Han,~A.~L.},
  Lastname,~C.~D.,
  Lastname,~E.~F.,
  {\it et al.}
  (2025)
  {\it Title of the paper},
  \href{https://doi.org/10.0000/example}{Journal Name~{\bf 16}, 2648}.
  \doi{10.0000/example}
\item
  Lastname,~G.~H.,
  {\bf Han,~A.~L.},
  Lastname,~I.~J.
  (2022)
  {\it Title of the earlier paper},
  \href{https://doi.org/10.0000/example2}{Journal Name~{\bf 19}(3), 190--196}.
  \doi{10.0000/example2}
% TODO: add further publications here, newest first
\end{etaremune}

% The presentation lists below are numbered independently of the publications,
% so no counter bookkeeping is needed -- just add new items at the top.
\section{Oral Presentations}
\secstartswithlist{}%
\begin{etaremune}
\item
  Venue or conference name, City, ST
  \hfill{}
  Month 2025
\item
  Venue or conference name, City, ST
  \hfill{}
  Month 2024
% TODO: add talks here, newest first
\end{etaremune}

\section{Poster Presentations}
\secstartswithlist{}%
\begin{etaremune}
\item
  Conference name, City, ST
  \hfill{}
  Month 2025
% TODO: add posters here, newest first
\end{etaremune}

\section{Honors and Awards}
% TODO: most recent first. Use the plain two-line idiom for anything without
% sub-details.
{\bf Award Name,} Granting Body \hfill {\bf 2024}\\
\\
{\bf Award Name,} Granting Body \hfill {\bf 2023--2024}\\
\\
{\bf Award Name,} Granting Body \hfill {\bf 2022}\\

% Added to improve page breaks
\vspace{-1em}

\section{Leadership and Entrepreneurship}
% TODO
{\bf Role,} Organization, \url{https://example.org} \hfill {\bf 2022--Present}
\begin{itemize}
\item[ ] TODO: what you built and its scale
\item[ ] TODO: outcome, coverage, or recognition
\end{itemize}

\section{Teaching and Mentoring}
% TODO
{\bf Role}, Institution
\vspace*{.05in}
\begin{itemize}
\item[ ] COURSE 000, Course Title \hfill {\bf Fall 2024}
\item[ ] COURSE 000, Course Title \hfill {\bf Spring 2024}
\end{itemize}

{\bf Mentees}
\begin{itemize}
\item[ ] Name, Institution \hfill {\bf 2024--Present}
\end{itemize}

\section{Service and Outreach}
% TODO
{\bf Professional Memberships}
\vspace*{.05in}
\begin{itemize}
\item[ ] Society Name \hfill {\bf 2023--Present}
\end{itemize}

{\bf Journal referee}
\vspace*{.05in}\\
\hspace*{1em}
TODO: comma-separated list of journals

{\bf Outreach}
\vspace*{.05in}
\begin{itemize}
\item[ ] Activity or venue \hfill {\bf 2024} \\
  \hspace*{1em} TODO: one line of detail
\end{itemize}

\section{Technical Skills}
% TODO
{\bf Languages---}%
TODO (e.g.\ Python, R, C/C++, Bash).
Markup: \LaTeX, Markdown, HTML/CSS.

{\bf Methods and tools---}%
TODO (e.g.\ statistical genetics, single-cell and bulk RNA-seq analysis,
machine learning, HPC/SLURM, Git).

{\bf Software---}%
TODO. Public code at \url{https://github.com/alhanster}.

\newpage{}

\section{References}
\vspace*{.05in}
% TODO: replace each block. Delete the ones you don't need.
\parbox{\textwidth}{%
{\bf Name,} Title, Institution \\
Street Address \\
City, ST 00000 \\
email: \href{mailto:name@example.edu}{name@example.edu}}
\par
\parbox{\textwidth}{%
{\bf Name,} Title, Institution \\
Street Address \\
City, ST 00000 \\
email: \href{mailto:name@example.edu}{name@example.edu}}
\par
\parbox{\textwidth}{%
{\bf Name,} Title, Institution \\
Street Address \\
City, ST 00000 \\
email: \href{mailto:name@example.edu}{name@example.edu}}

\end{resume}
\end{document}

% Local Variables:
% mode: latex
% End:
